% !PRINTED 182
\weberpage{182}

%% [running head, not body text -- commented out 2026-08-07] Page furniture (running
%% head, folio, signature) is excluded from this edition by standing policy.
%% \noindent 182 \hfill \textit{Sechster Abschnitt.} \hfill \S.\,47.

%<<CONT
haben, rational durch $\Theta$ und die $A_k$ ausdrücken, und wir setzen
\weq{4}{x_i = \varphi_i(\Theta, A_1, A_2, \ldots).}

Auch diese Ausdrücke werden, wenn über die Function $\Theta$ verfügt ist, für die Normalform feste Gestalten annehmen.

Nun kann aber auch der Fall vorkommen, dass die Invarianten nicht in der Normalform, sondern in einer beliebigen anderen Form, die wir die allgemeine Form nennen wollen, gegeben sind. Dann wird es sich darum handeln, die Substitution $L$ zu finden, durch die die allgemeine Form in eine von vornherein als möglich erkannte Normalform transformirt wird. Diese Aufgabe ist, wie wir nun sehen wollen, keine andere, als das für die Normalform gestellte Formenproblem selbst.

Nehmen wir an, die Invarianten in der allgemeinen Form, als Functionen von $y$, seien $B_1, B_2, \ldots$, so dass durch die Substitution (2) die Identitäten
\weq{5}{A_1 = B_1, \ A_2 = B_2, \ldots}
hergestellt werden. Es ist die Aufgabe, wenn die Functionen $A_i, B_i$ der Form nach gegeben sind, die Substitution $L$ zu bestimmen, die die Gleichungen (5) zu identischen macht. Diese Aufgabe ist gelöst, wenn wir die Gleichung (3) für ein passend gewähltes specielles Werthsystem der $A_i$ als gelöst voraussetzen. Um dies nachzuweisen, bilden wir die vollständigen Differentiale der Gleichungen (3) und (4):
\[
0 = \Phi'(\Theta)\,d\Theta + \sum^s \Phi'(A_s)\,dA_s
\]
\[
dx_i = \varphi_i'(\Theta)\,d\Theta + \sum^s \varphi_i'(A_s)\,dA_s,
\]
worin $\Phi'(\Theta)$, $\Phi'(A_s)$, $\varphi_i'(\Theta)$, $\varphi_i'(A_s)$ die partiellen Ableitungen bedeuten. Eliminiren wir $d\Theta$, so folgt:
\weq{6}{dx_i = \sum^s \frac{\varphi_i'(A_s)\,\Phi'(\Theta) - \Phi'(A_s)\,\varphi_i'(\Theta)}{\Phi'(\Theta)}\, dA_s.}

Nun ist in Folge der Gleichungen (5):
\weq{7}{dA_s = \sum^k B_s'(y_k)\, dy_k,}
und wenn wir also
\weq{8}{x_i = \alpha_{1,i} y_1 + \cdots + \alpha_{n,i} y_n}
\weq{9}{dx_i = \alpha_{1,i}\, dy_1 + \cdots + \alpha_{n,i}\, dy_n}
%CONT>>
